\begin{itemize}
    \color{red}
    \item key algorithm for track reconstruction
    \item one of the most expensive algorithms in terms of CPU time
    \item needs to be fast and efficient
    \item ACTS comes with a combinatorial Kalman filter (CKF) implementation
    \begin{itemize}
        \item new alternatives: GNN or transformer networks
    \end{itemize}
    \item relies on starting parameters which can be obtained from triplet seeds
    \item CKF has been rewritten during the work on this thesis; more flexible, more efficient, more maintainable
\end{itemize}

\section{The combinatorial Kalman filter}

\begin{itemize}
    \color{red}
    \item how does the algorithm work?
    \begin{itemize}
        \item relies on transport of track parameters and their covariance
        \item source link accessor
        \item how are the track parameters updated? shared with KF
        \item calibration
        \item measurement selection
        \item branch stopper
    \end{itemize}
    \item optimal measurement discrimination power
    \item can provide fully fitted tracks after smoothing pass
    \item exercised in TrackFindingAlgorithm in Examples
\end{itemize}

\section{CKF as a component}

\begin{itemize}
    \color{red}
    \item CKF used to be a monolithic algorithm which does finding, fitting, smoothing, and extrapolation
    \item Code was hard to read and follow, bugs resulted in low efficiency
    \item CKF has been reduced to be a track finding component; only does the finding
\end{itemize}

\section{Two-way track finding}

\begin{itemize}
    \color{red}
    \item Seeds are not guaranteed to start on the first measurement
    \item Track finding needs to be able to look for measurements in both directions
    \item After the CKF has been turned into a component, this was implemented by running the CKF twice
\end{itemize}

\section{Branch stopper}

\begin{itemize}
    \color{red}
    \item can safe time by stopping branches early if they are not promising
    \item allows to continue, stop and drop, and stop and keep
\end{itemize}

\section{Track container for branches}

\begin{itemize}
    \color{red}
    \item branches during CKF ressemble tracks; have measurement count, hole count, outlier count, and chi2
    \item branches are now stored as tracks in a track container
    \item branch stopper can observe these tracks and can accumulate statistics
    \item extensible for experiments through dynamic columns
\end{itemize}

\section{Branch finalization}

\begin{itemize}
    \color{red}
    \item branch finalization used to be done in multiple places
    \item now done with a single function which can be called from different places
\end{itemize}

\section{Seed deduplication}

\begin{itemize}
    \color{red}
    \item triplet seeding results in duplicates
    \item this is partially by design to gain redundancy and to mitigate dead modules
    \item deduplication is done during track finding
    \item seeds that have been fully covered by an already found track will be skipped
\end{itemize}

\section{Stay-on-seed strategy}

\begin{itemize}
    \color{red}
    \item CKF only uses starting parameters
    \item given a seed it makes sense to stay on the selected measurements
    \item can be implemented via source link accessor
\end{itemize}

\section{Initial covariance}

\begin{itemize}
    \color{red}
    \item for unbiased fitting initial covariance needs to be inflated
    \item during finding the covariance also dictates the search range; with branching it cannot be too large
    \item cite frühwrith
\end{itemize}

\section{Shortcomings}

\begin{itemize}
    \color{red}
    \item generic measurement selector
    \item extended bounds for finding
    \item edge holes
    \item multi-surface measurement selection
\end{itemize}
